\documentclass[../main.tex]{subfiles}
\begin{document}

\chapter{イントロダクション}
\lessoninfo{
  video={Lesson 1，動画 1--19},
  textbook={H\&P \cite{hennessy2017} §1.1--1.6},
  keywords={コンピュータアーキテクチャ，ムーアの法則，メモリウォール，動的電力，静的電力，製造コスト，歩留まり},
}

この章では，コンピュータアーキテクチャとは何かを定義し，それが独立した
学問分野として必要とされる理由を述べ，以降のすべての設計判断を制約する
技術トレンドと3つの量---性能・電力・コスト---を導入する．

\section{コンピュータアーキテクチャとは何か}

\term[こんぴゅーたあーきてくちゃ]{コンピュータアーキテクチャ}[computer architecture]
とは，利用可能な製造技術を用いて，目的に適したコンピュータを設計する
学問分野である．この定義には強調すべき点が2つある．\source{Lesson 1，動画
2--3；H\&P §1.1--1.3．}

\begin{itemize}
  \item \emph{目的に適していること}．携帯電話，ノート PC，データセンタの
    サーバはいずれもコンピュータだが，最適化の対象が異なる．一方では電池
    寿命・大きさ・重さ，他方では生の性能やコスト当たりのスループットで
    ある．したがってアーキテクチャとは唯一の最良設計を求めることではなく，
    \emph{トレードオフ}を扱うことである．
  \item \emph{利用可能な技術}．アーキテクトはトランジスタそのものを設計
    しないが，現在のプロセスで何が実現できるか---1 チップに載る
    トランジスタ数，スイッチング速度，消費電力---を知っていなければ
    ならない．それらが達成可能な範囲を定めるからである．\margin{規模の目安:
    Intel 4004（1971 年）のトランジスタ数は 2300，NVIDIA H100（2022 年）は
    800 億である．ベンダ公表値．}
\end{itemize}

\subsection{なぜ必要か}

コンピュータアーキテクチャの目標は2つある．\margin{講義では「性能の向上」
と「能力の向上」とまとめられている．}
\begin{enumerate}
  \item \textbf{性能の向上}（広い意味で）: 実行速度，電池寿命，大きさ，
    重さ，エネルギー効率．
  \item \textbf{能力の向上}: これまでできなかったことを可能にする．例えば
    3D グラフィックス，デバッグのハードウェア支援，セキュリティ機能．
\end{enumerate}

もう1つの理由はタイミングである．今日設計するプロセッサが製造されるのは
数年後なので，アーキテクトは現在のパラメータではなく，その時点で技術が
どこにあるかを\emph{予測}して設計しなければならない．\tip{後の章で
過剰設計に見える判断があったら，それがどの技術トレンドを見越したもの
かを考えるとよい．}

\section{技術トレンド}

\subsection{ムーアの法則}

\term[むーあのほうそく]{ムーアの法則}[Moore's law]とは，一定面積のチップに
載せられるトランジスタ数が約 18〜24 か月ごとに倍増するという経験則で
ある．\source{\cite{moore1965}；Lesson 1，動画 6；H\&P §1.4．}物理法則
ではなくトレンドだが，数十年にわたって成り立ち，業界全体の期待を形作って
きた．講義では，同じ 18〜24 か月の倍増周期で次の帰結が導かれる．
\caution{ムーア自身は 1965 年に1年で倍増と述べ，1975 年に2年へ改めた．
よく引かれる「18 か月」は David House の数字であり，ムーアのものではない．}
\begin{itemize}
  \item プロセッサの\emph{速度}が倍になる（トランジスタが増えれば並列性を
    高められ，ゲートは小さく速くなる）．
  \item \emph{演算当たりのエネルギー}が半分になる．
  \item \emph{メモリ容量}が倍になる．
\end{itemize}

\begin{example}
  速度が2年ごとに倍になるなら，10 年で $2^{5} = 32$ 倍の向上が期待される．
  逆に言えば，4年かけて従来品の2倍しか速くないプロセッサを作った設計
  チームは，トレンドに追いついたにすぎない．
\end{example}

\subsection{メモリウォール}

すべての構成要素が同じ曲線に従うわけではない．プロセッサ性能が 18〜24
か月で倍増する一方，メモリの\emph{レイテンシ}は同じ期間に約 $1.1$ 倍しか
改善しなかった．\source{\cite{wulf1995}；Lesson 1，動画 8；H\&P §2.1．}
その結果，プロセッサがデータを消費できる速さとメモリがそれを供給できる
速さの差は世代ごとに広がる．この差を
\term[めもりうぉーる]{メモリウォール}[memory wall]と呼び，キャッシュ
（第 12〜14 章）が重要である理由はここにある．

\begin{figure}[htbp]
  \centering
  % Processor vs. memory performance trend (schematic, log scale).
\begin{tikzpicture}[x=1cm,y=1cm,>={Stealth[length=2mm]}, font=\small\sffamily]
  \draw[->] (0,0) -- (6.4,0) node[below left] {\iflangja{年}{year}};
  \draw[->] (0,0) -- (0,4.2) node[left] {\iflangja{性能（対数）}{performance (log)}};
  \draw[thick,ln-accent] (0.3,0.3) -- (6,3.9) node[pos=0.75,above left] {\iflangja{プロセッサ}{processor}};
  \draw[thick,ln-caution] (0.3,0.3) -- (6,1.2) node[pos=0.8,below] {\iflangja{メモリ（遅延）}{memory (latency)}};
  \draw[<->,ln-muted] (5.6,1.15) -- (5.6,3.65) node[midway,right,font=\footnotesize] {\iflangja{差}{gap}};
\end{tikzpicture}

  \caption{プロセッサ性能とメモリレイテンシは異なる速さで改善する．その差が
    メモリウォールである．}
  \label{fig:memory-wall}
\end{figure}

\begin{keypoint}
  メモリ容量はムーアの法則に追随するが，メモリの\emph{レイテンシ}は追随
  しない．本講義のメモリシステム設計の大半は，このレイテンシを隠すために
  存在する．\margin{Skylake では L1 ヒットが 4 サイクル，L2 ヒットが 12
  サイクル，DRAM アクセスは 60〜\SI{90}{\nano\second}，すなわち
  \SI{3}{\giga\hertz} では 200〜300 サイクルに相当する．Intel 最適化
  マニュアル；7-cpu.com．}
\end{keypoint}

\section{速度・電力・コスト}

長い間「より良い」とは単に「より速い」を意味した．今日のアーキテクトは
3つの量のバランスを取らなければならない．\source{Lesson 1，動画 9--10．}
\begin{description}
  \item[性能] 想定するワークロードをどれだけ速く実行できるか．
  \item[電力] どれだけのエネルギーを消費するか．携帯機器では電池寿命を，
    サーバでは冷却コストを決める．\margin{電力予算: 携帯 SoC は
    2〜\SI{5}{\watt}，デスクトップ CPU は 65〜\SI{125}{\watt}，サーバ CPU は
    最大 \SI{400}{\watt}．ベンダ公表値．}
  \item[コスト] 製造にいくらかかるか．
\end{description}
どれが支配的かは製品による．携帯電話なら重さと電池寿命，組込み
コントローラならコスト，スーパーコンピュータなら性能である．
\caution{ワークロードと電力・コスト予算を伴わない「より速く」は，
意味のある設計目標ではない．}

\section{消費電力}

プロセッサが消費する電力は2つの成分からなる．\source{Lesson 1，動画
11--13；H\&P §1.5．}

\subsection{動的電力（アクティブ電力）}

\term[どうてきでんりょく]{動的電力}[dynamic (active) power]は，信号が
切り替わるたびに配線やゲートの容量を充放電することで消費される．
\begin{equation}
  P_{\mathrm{dyn}} = \tfrac{1}{2}\, C\, V^{2}\, f\, \alpha
  \label{eq:active-power}
\end{equation}
ここで $C$ はスイッチングされる総容量，$V$ は電源電圧，$f$ はクロック
周波数，$\alpha$ は\emph{活性化率}，すなわち各サイクルで実際に切り替わる
容量の割合である．\margin{$C$ はチップ面積とトランジスタ数とともに増え，
$\alpha$ はワークロードと，どれだけの論理が遊んでいるかに依存する．
H\&P は $f\alpha$ の積を「スイッチする周波数」と書く．}

\cref{eq:active-power} から直ちに2つのことが分かる．
\begin{itemize}
  \item 電力は電圧の\emph{2乗}に比例するので，$V$ を下げるのが最も効果的な
    手段である．ただし回路が維持できる最大周波数も $V$ とともに下がるため，
    実際には $V$ と $f$ を同時に下げることになり，電力はおよそ $V^{3}$ に
    比例して減る．
  \item 電圧一定で周波数を倍にすると電力も倍になる．同じ高速化を
    トランジスタ数の倍増（並列ハードウェアの追加）で得る場合，電力の
    増加は $C$ の増分だけである．\sidenote{2005 年ごろまでは，新しいプロセス
    ごとにトランジスタとともに $V$ も縮小され，同じ電力のままチップを高速化
    できた．リーク電流により $V$ の縮小は \SI{1}{\volt} 付近で止まり，クロック
    周波数は 3〜\SI{4}{\giga\hertz} で頭打ちとなって，増えたトランジスタは
    コア数の増加（第 18 章）に向けられた．}
\end{itemize}

\begin{example}
  $C$，$f$，$\alpha$ を固定したまま電源電圧を \SI{15}{\percent} 下げると，
  動的電力は元の $0.85^{2} \approx 0.72$ 倍，すなわち約 \SI{28}{\percent}
  減る．
\end{example}

\subsection{静的電力}

\term[せいてきでんりょく]{静的電力}[static (leakage) power]は，信号が
切り替わらなくても消費される．トランジスタは完全なスイッチではなく，
オフでも電流が漏れるからである．リーク電流は，電源電圧---それに伴って
トランジスタのしきい値電圧---を下げるほど大きくなる．したがって
トレードオフがある．$V$ を下げると動的電力は減るが静的電力は増え，
合計 $P = P_{\mathrm{dyn}} + P_{\mathrm{static}}$ はある中間の電圧で
最小になる（\cref{fig:power}）．

\begin{marginfigure}
  \resizebox{\linewidth}{!}{% Active vs. static power as a function of supply voltage (schematic).
\begin{tikzpicture}[x=1cm,y=1cm,>={Stealth[length=2mm]}, font=\small\sffamily]
  \draw[->] (0,0) -- (6.2,0) node[below left] {\iflangja{電源電圧 $V$}{supply voltage $V$}};
  \draw[->] (0,0) -- (0,4.2) node[left] {\iflangja{電力}{power}};
  % active ~ V^2
  \draw[thick,ln-accent] plot[domain=0.6:5.6,samples=40] (\x,{0.11*\x*\x});
  \node[ln-accent,anchor=west] at (4.3,3.4) {\iflangja{動的電力 $\propto V^2$}{active $\propto V^2$}};
  % static: grows as V drops (leakage at low threshold)
  \draw[thick,ln-caution] plot[domain=0.6:5.6,samples=40] (\x,{1.6/\x});
  \node[ln-caution,anchor=west] at (0.9,2.9) {\iflangja{静的電力}{static (leakage)}};
  % total
  \draw[thick,dashed,ln-muted] plot[domain=0.6:5.6,samples=40] (\x,{0.11*\x*\x+1.6/\x});
  \node[ln-muted,anchor=south west] at (4.0,1.1) {\iflangja{合計}{total}};
  \draw[dotted] (1.94,0) -- (1.94,{0.11*1.94*1.94+1.6/1.94});
  \node[below,font=\footnotesize] at (1.94,0) {$V_{\mathrm{opt}}$};
\end{tikzpicture}
}
  \caption{電源電圧を変えると，動的電力と静的電力は逆方向に動く．}
  \label{fig:power}
\end{marginfigure}

\begin{keypoint}[電力]
  動的電力 $\propto C V^{2} f \alpha$ は低電圧を好み，静的電力は低電圧を
  罰する．現代の設計は両者の和の最小点付近にあり，さらなる改善は $V$ より
  も $C$，$\alpha$，$f$ の削減から得ている．
\end{keypoint}

\section{製造コスト}

チップは円形のシリコン\term[うぇは]{ウェハ}[wafer]上に製造され，各ウェハは
多数の同一の長方形の\term[だい]{ダイ}[die]に切り分けられる．ウェハ1枚を
処理するコストはほぼ一定なので，ダイ1個のコストは次式で与えられる．
\source{Lesson 1，動画 15--17；H\&P §1.6．}
\begin{equation}
  \text{ダイ当たりコスト} \;=\;
  \frac{\text{ウェハ当たりコスト}}
       {\text{ウェハ当たりダイ数} \times \text{歩留まり}}
  \label{eq:die-cost}
\end{equation}

\subsection{歩留まり}

すべてのダイが動作するわけではない．ウェハ上には欠陥が散在し，欠陥を
含むダイは通常廃棄される．\term[ぶどまり]{歩留まり}[yield]とは，機能する
ダイの割合である．欠陥はランダムな位置に生じるため，大きなダイほど欠陥を
含みやすく，歩留まりはダイ面積が増えるにつれて\emph{低下}する．
\margin{H\&P §1.6 には経験的モデル
$\text{歩留まり} \approx 1/(1 + \text{欠陥密度}\times\text{面積})^{N}$
が示されている．成熟したプロセスの欠陥密度は \si{\square\centi\metre} 当たり
\num{0.1} 程度である．}

\subsection{帰結}

\cref{eq:die-cost} と歩留まりの振る舞いを合わせると，この節の中心的な
結論が得られる．ダイ面積を半分にすると，ダイ当たりコストは半分\emph{以下}
になる．ウェハ当たりのダイ数が倍になり，\emph{かつ}歩留まりも上がるから
である．\tip{ウェハ当たりダイ数 $\approx \pi (d/2)^{2}/A - \pi d/\sqrt{2A}$．
第2項はウェハ端の欠けたダイを除く分である．H\&P §1.6．}ムーアの法則と
合わせて読むと，アーキテクトには技術世代ごとに2つの選択肢がある．
\begin{itemize}
  \item ダイ面積を一定に保ち，2倍のトランジスタを詰め込む---ほぼ同じ
    コストでより高機能なチップ．\margin{\SI{300}{\milli\metre} ウェハの面積は
    約 \SI{707}{\square\centi\metre}，リソグラフィの1フィールドで描ける最大の
    ダイは $26 \times \SI{33}{\milli\metre} = \SI{858}{\square\milli\metre}$
    であり，NVIDIA H100 はその \SI{814}{\square\milli\metre} を使う．}
  \item 設計を一定に保ち，面積を半分に縮小する---同じチップを半分よりずっと
    安いコストで．
\end{itemize}

\begin{example}
  ウェハが \$5000 で，ある設計のダイが 100 個取れ，歩留まりが
  \SI{50}{\percent} だとする．動作するダイ1個のコストは
  $5000/(100\times 0.5) = \$100$ である．縮小によりウェハ当たり 200 個
  取れ，歩留まりが \SI{70}{\percent} に上がると，コストは
  $5000/(200\times 0.7) \approx \$36$ と半分以下になる．
\end{example}

\begin{keypoint}[章のまとめ]
  アーキテクチャとは，技術トレンド（ムーアの法則）とその例外（メモリ
  ウォール）に駆動されながら，性能・電力・コストという3つの制約の下で
  トレードオフを設計することである．次章では「性能」を厳密に定義する．
\end{keypoint}

\end{document}
